\section{User Interaction and Application Behavior}
The final implementation of the system successfully fulfills the primary 
functional objectives defined at the beginning of the project. 
% The application enables users to search for a geographic region, 
% retrieve its associated local cuisines, and obtain restaurant 
% recommendations based on a selected cuisine and location. 
The system integrates a Spring Boot backend with a Next.js frontend, 
ensuring seamless communication through RESTful APIs.
In the main page, users are able to explore multiple geographic regions.
A list of available regions are listed, and retrieve the cuisines associated with a specific region. 
The regions correspond to administrative divisions, such as states or prefectures, depending on the respective country.

\input{myFigures/FiguresCommand.tex}
\regionsFigure

Once a region is selected, users are able to see the popular cuisines that represent the area.

\cuisinesFigure

When a user selects a specific cuisine, the user is directed to the detailed 
cuisine page. Details about the cuisine, such as its description
and representative characteristics, are displayed to the user 
in a structured format. Reviews associated with the cuisine are also loaded
on the page, allowing users to view previously submitted reviews. 
The interface also provides users with the ability to submit new reviews 
through an input form. Upon submission, the review is sent to the backend 
through an API request and stored in the database. The updated review 
list is then dynamically rendered without requiring a full page reload, 
providing a seamless user experience. 
In addition to browsing cuisine information, users have the option 
to search for restaurants near their selected geographic area. 
When this button is clicked, the application opens a new browser tab 
that redirects the user to a Google Maps search page. 
The search query is automatically generated based on the selected 
cuisine and the user's location, allowing users to view nearby 
restaurants relevant to their chosen cuisine. 

By applying Google Maps search for location-based search, 
the system provides users with real-world, up-to-date restaurant 
results without requiring the application to maintain its own 
mapping infrastructure.

\detailedCuisinesFigure

The frontend interface dynamically updates content based on user actions, 
ensuring that navigation between regions, cuisines, reviews, and 
restaurant recommendations occurs smoothly. This interactive behavior 
demonstrates the successful integration of the frontend and backend 
components, as well as the correct handling of asynchronous API requests.

Overall, the system provides an intuitive workflow in which users 
can explore regions, learn about local cuisines, interact through 
reviews, and obtain nearby restaurant recommendations within a 
single application.



